% !TeX root = ../literature-review.tex
\section{Literature Review}
\label{sec:literature-review}

This chapter reviews research on Text-to-SQL systems, 
particularly in terms of correctness, 
recovering from errors and evaluating query trajectories.  

\subsection{Text-to-SQL Systems and Their Evaluation}
\label{sec:literature-foundations}

\subsubsection{From Natural Language To SQL}

The primary task of Text-to-SQL systems is to translate natural-language 
questions into SQL queries that answer the user's request.
These systems enable users to interact with databases in natural language 
and can be beneficial to non-expert users 
who may not be familiar with SQL \citep{shi2025survey}.
It is, however, necessary to use the correct tables and columns 
in the generated query to answer the question.
Merely outputting a valid SQL query that the database can execute is therefore insufficient.

\citet{shi2025survey} describe the progression of Text-to-SQL systems from 
rule and template-based to learning approaches.  
Systems leveraging LLMs benefit from pre-training  
and the ability to learn from in-context examples. 
\citet{shi2025survey} classify these approaches into prompting
and fine-tuning. Fine-tuning further trains models on task-related data.

Prompt design can affect the correctness of generated SQL.
\citet{gao2024benchmark} examined question representation,
example selection and prompt organisation, showing how these choices
affected Text-to-SQL performance.

Schema linking refers to associating the relevant database entities with the question.
In their survey, \citet{shi2025survey} describe it as a step 
to build the context used to prompt the model to produce an SQL query.
This includes the table and column names and their relationships.
Depending on the question, producing the correct query 
may require information from several related tables.

\subsubsection{Benchmarks and Execution Feedback}

Existing Text-to-SQL benchmarks evaluate different aspects of the systems.
For example, Spider introduces the challenge of unfamiliar database schemas at test time \citep{yu2018spider}.
BIRD additionally tests the use of database contents and external knowledge, 
as well as the efficiency of SQL queries \citep{li2023bird}.
Hence, differences in performance on these benchmarks may reflect the different challenges they present.

Benchmark evaluation also depends on how correctness is measured.
\citet{gao2024benchmark} describe execution accuracy as agreement
between the results of generated and reference queries.
A query can execute successfully and still return different data.

The survey by \citet{shi2025survey} describes SQL correction methods
that use returned results or execution errors to check and revise a query.
The final result may therefore follow several SQL attempts.
Checking that result shows whether it matches the expected answer.
Examining the attempts also shows which errors occurred and how
the queries changed after database feedback.